% ============================================================
%  第五章：实验评估 —— 独立编译版本
%  直接上传到 Overleaf，用 XeLaTeX 编译
% ============================================================
\documentclass[12pt,a4paper]{article}
\usepackage{geometry}
\usepackage{booktabs}
\usepackage{hyperref}
\usepackage{tocloft}
\usepackage{multirow}
\usepackage{array}
\usepackage{amsmath}
\usepackage{ctex}
\usepackage{fancyhdr}
\geometry{margin=2.5cm}
\hypersetup{colorlinks=true, linkcolor=black, urlcolor=black}

\usepackage{xurl}
\urlstyle{tt}
\usepackage{graphicx}
\usepackage{longtable,tabularx,array,multirow}
\usepackage{setspace}
\usepackage{xcolor}
\usepackage{amssymb,bm}
\usepackage{caption}
\usepackage{float}
\usepackage{tikz}
\usetikzlibrary{arrows.meta,positioning,fit,calc,shapes.geometric,shapes.symbols,matrix,backgrounds,decorations.pathreplacing}

\usepackage{enumitem}
\setlist[itemize]{leftmargin=2em,itemsep=0.25em,topsep=0.25em}
\setlist[enumerate]{leftmargin=2em,itemsep=0.25em,topsep=0.25em}

\definecolor{DeepBlue}{HTML}{163A5F}
\definecolor{TechBlue}{HTML}{2563EB}
\definecolor{Mint}{HTML}{DDF7EF}
\definecolor{SoftBlue}{HTML}{E8F1FF}
\definecolor{SoftOrange}{HTML}{FFF3E0}
\definecolor{SoftRed}{HTML}{FFE7E7}
\definecolor{SoftPurple}{HTML}{F0E8FF}
\definecolor{LineGray}{HTML}{64748B}
\definecolor{DarkText}{HTML}{0F172A}

\newcommand{\projectname}{基于时序可信动态图谱的新能源锂电池产业链黑天鹅风险推理与预警系统}
\newcommand{\projectshort}{TIGER-Risk}
\newcommand{\term}[1]{\textbf{#1}}

\usepackage{listings}
\renewcommand{\lstlistingname}{代码}
\lstdefinestyle{projectcode}{
  language=JavaScript,
  basicstyle=\ttfamily\scriptsize,
  numbers=left,
  numberstyle=\tiny\color{LineGray},
  stepnumber=1,
  numbersep=6pt,
  frame=single,
  rulecolor=\color{LineGray},
  backgroundcolor=\color{SoftBlue},
  breaklines=true,
  breakatwhitespace=false,
  columns=fullflexible,
  keepspaces=true,
  showstringspaces=false,
  tabsize=2,
  captionpos=b,
  xleftmargin=1em,
  xrightmargin=1em,
  aboveskip=0.8em,
  belowskip=0.8em
}

\pagestyle{fancy}
\fancyhf{}
\fancyhead[C]{\small 基于时序可信动态图谱的新能源锂电池产业链黑天鹅风险推理与预警系统}
\renewcommand{\headrulewidth}{0.6pt}
\setlength{\headheight}{15pt}
\setlength{\headsep}{18pt}
\fancypagestyle{plain}{
    \fancyhf{}
    \fancyhead[C]{\small 基于时序可信动态图谱的新能源锂电池产业链黑天鹅风险推理与预警系统}
    \renewcommand{\headrulewidth}{0.6pt}
}

\begin{document}

% ====================== 5 实验评估 ======================
\clearpage
\section{实验评估}

本章围绕\projectshort（以下简称\projectshort）系统进行全面的实验评估。与分类算法或回归模型的单一指标评测不同，\projectshort 是一个面向新能源锂电池产业链黑天鹅风险推理的端到端智能系统，包含多源数据接入、事件抽取、时序动态图谱构建、Trust Graph 可信度评估、图增强风险传播推理和大语言模型研判报告六个核心模块。因此，本章的实验设计遵循系统级评估范式：\term{提出假设$\to$设计实验$\to$验证假设$\to$分析原因$\to$得出结论}，而非简单堆叠 Accuracy/Precision/Recall 指标。

\begin{figure}[H]
\centering
\resizebox{\linewidth}{!}{
\begin{tikzpicture}[
    font=\small,
    box/.style={rectangle, rounded corners=3pt, draw=DeepBlue, thick, fill=SoftBlue, align=center, minimum width=2.8cm, minimum height=0.85cm},
    core/.style={rectangle, rounded corners=3pt, draw=TechBlue, thick, fill=Mint, align=center, minimum width=2.8cm, minimum height=0.85cm},
    arr/.style={-{Latex[length=2mm]}, thick, draw=LineGray}
]
\node[box] (h) {实验假设\\H1--H5};
\node[core, right=0.6cm of h] (p) {实验协议\\公平性+重复性};
\node[core, right=0.6cm of p] (d) {数据集\\构建与标注};
\node[box, below=0.7cm of d] (e1) {实验1:\\事件抽取};
\node[box, left=0.6cm of e1] (e2) {实验4:\\效率分析};
\node[box, right=0.6cm of e1] (e3) {实验2:\\图谱检索};
\node[box, below=0.7cm of e1] (e4) {实验3:\\风险预测};
\node[core, below=0.7cm of e4] (cs) {Case Study};
\node[core, below=0.7cm of cs] (ab) {Ablation\\+参数敏感性};
\node[box, below=0.7cm of ab] (di) {讨论与\\效度威胁};
\draw[arr] (h) -- (p);
\draw[arr] (p) -- (d);
\draw[arr] (d.south) -- ++(0,-0.25) -| (e1.north);
\draw[arr] (d.south) -- ++(0,-0.25) -| (e2.north);
\draw[arr] (d.south) -- ++(0,-0.25) -| (e3.north);
\draw[arr] (d.south) -- ++(0.7,0) |- (e4.north);
\draw[arr] (e1.south) -- ++(0,-0.1) -- ++(1.5,0) |- (cs.east);
\draw[arr] (e2.south) -- ++(0,-0.1) -- ++(-1.5,0) |- (cs.west);
\draw[arr] (e3.south) -- ++(0,-0.1) -- ++(1.5,0) |- (cs.east);
\draw[arr] (e4.south) -- (cs.north);
\draw[arr] (cs.south) -- (ab.north);
\draw[arr] (ab.south) -- (di.north);
\end{tikzpicture}
}
\caption{实验评估总体框架}
\label{fig:exp-framework}
\end{figure}

本章整体组织如下：第~\ref{sec:exp-framework}~节建立四层评价体系并提出实验假设；第~\ref{sec:exp-protocol}~节定义实验协议与公平性保障；第~\ref{sec:exp-env}~节说明实验环境；第~\ref{sec:exp-dataset}~节介绍数据集构建与 Ground Truth 标注；第~\ref{sec:exp-baseline}~节定义基线方法；第~\ref{sec:exp-metrics}~节列举全部评价指标；第~\ref{sec:exp-extraction}~至~\ref{sec:exp-efficiency}~节依次报告四个核心实验；第~\ref{sec:exp-casestudy}~节给出完整案例推演；第~\ref{sec:exp-ablation}~节进行消融实验与参数敏感性分析；第~\ref{sec:exp-discussion}~节讨论实验发现；第~\ref{sec:exp-threats}~节分析效度威胁。

\subsection{评价体系与实验假设}
\label{sec:exp-framework}

\subsubsection{四层评价体系设计}

由于\projectshort 系统包含多个串联模块，单一指标（如分类 Accuracy）无法全面评价系统质量。本文设计四层评价体系，分别覆盖系统的不同能力维度。

\begin{enumerate}
    \item \term{语义理解层（L1）}：评价系统从非结构化文本中提取结构化风险事件的能力。对应实验为事件抽取实验（第~\ref{sec:exp-extraction}~节）。
    \item \term{知识检索层（L2）}：评价系统从时序动态图谱中检索相关事件和影响路径的能力。对应实验为图谱检索实验（第~\ref{sec:exp-retrieval}~节）。
    \item \term{风险推理层（L3）}：评价系统的风险等级预测、幻觉抑制和可信度判断能力。对应实验为风险预测实验（第~\ref{sec:exp-prediction}~节）。
    \item \term{工程效能层（L4）}：评价系统的响应速度、吞吐量和资源开销。对应实验为效率分析实验（第~\ref{sec:exp-efficiency}~节）。
\end{enumerate}

\subsubsection{实验假设}

基于\projectshort 系统的核心技术主张，本文提出以下五个实验假设：

\begin{itemize}
    \item \term{H1（时序增强检索假设）}：时序动态图谱能够显著提升产业链风险事件的检索能力，即：
    \begin{equation}
    \text{Hit@K}_{\text{Temporal}} > \text{Hit@K}_{\text{Static}}, \quad \text{MRR}_{\text{Temporal}} > \text{MRR}_{\text{Static}}.
    \end{equation}
    该假设的直观逻辑在于：带时间状态的图谱能够过滤过期关系、优先召回有效事件，从而提高检索精度。

    \item \term{H2（可信增强预测假设）}：Trust Graph 多源可信度聚合能够提升风险等级预测的准确率，即：
    \begin{equation}
    \text{Accuracy}_{\text{Trust}} > \text{Accuracy}_{\text{w/o Trust}}, \quad \text{F1}_{\text{Trust}} > \text{F1}_{\text{w/o Trust}}.
    \end{equation}

    \item \term{H3（图增强幻觉抑制假设）}：将结构化图谱路径作为大语言模型的输入约束，能够显著降低生成报告中的事实性幻觉。本文使用\textbf{证据对齐精度}（Evidence Alignment Precision, EAP）衡量生成内容与输入证据之间的一致性：
    \begin{equation}
    \text{EAP} = \frac{\text{报告中与证据一致的断言数}}{\text{报告中总断言数}}.
    \end{equation}
    假设形式化为：$\text{EAP}_{\text{Graph-enhanced}} > \text{EAP}_{\text{Plain LLM}}$。

    \item \term{H4（效率可控假设）}：引入时序动态图谱和 Trust Graph 带来的额外计算开销在可接受范围内。定义效率可接受条件为：
    \begin{equation}
    \frac{\text{Latency}_{\text{TIGER}}}{\text{Latency}_{\text{Base}}} < 1.25 \quad \text{且} \quad \frac{\text{Accuracy}_{\text{TIGER}}}{\text{Accuracy}_{\text{Base}}} > 1.10.
    \end{equation}

    \item \term{H5（协同增益假设）}：时序动态图谱、Trust Graph 和图增强 LLM 三者协同工作的效果优于各组件单独增益之和，即：
    \begin{equation}
    \Delta_{\text{Full}} > \Delta_{\text{Temporal}} + \Delta_{\text{Trust}} + \Delta_{\text{Graph-LLM}}.
    \end{equation}
    该假设通过消融实验（第~\ref{sec:exp-ablation}~节）验证。
\end{itemize}

表~\ref{tab:hypothesis-mapping}~总结了各实验与假设之间的对应关系。

\begin{table}[H]
\centering
\footnotesize
\caption{实验-假设映射关系}
\label{tab:hypothesis-mapping}
\begin{tabularx}{0.95\linewidth}{c c c X}
\toprule
\textbf{实验} & \textbf{验证假设} & \textbf{关键指标} & \textbf{说明} \\
\midrule
Event Extraction & H1（支撑证据） & 字段完整率、F1 & 事件抽取质量直接影响后续检索；低质量抽取会导致检索基准不公平，因此需先确认各系统抽取能力相当 \\
Graph Retrieval & H1（主验证） & Hit@K、MRR、Recall@K & 直接比较静态图谱与动态图谱的检索能力 \\
Risk Prediction & H2、H3 & Accuracy、EAP、F1 & 同时验证可信度对预测准确率的提升和图结构对幻觉的抑制 \\
Efficiency & H4 & Latency P50/P95/P99、TPS & 验证引入动态图与可信机制的代价是否可接受 \\
Ablation Study & H5 & 全部指标按组件移除 & 验证各组件贡献及协同效应 \\
\bottomrule
\end{tabularx}
\end{table}

\subsection{实验协议}
\label{sec:exp-protocol}

为保证实验的公平性、可重复性和统计可靠性，本章遵循以下实验协议。

\subsubsection{公平性原则}

所有对比方法在以下维度保持严格统一：
\begin{itemize}
    \item \term{大语言模型}：统一使用 DeepSeek V4 Pro，Temperature = 0.1，Top-p = 0.9，Max Tokens = 4096；
    \item \term{文本嵌入模型}：统一使用 \texttt{text-embedding-3-large}，维度 3072；
    \item \term{随机种子}：统一设置 Python Seed = 42，NumPy Seed = 42，PyTorch Seed = 42；
    \item \term{知识图谱存储}：Neo4j Community 5.20.0，使用同一快照的产业链种子图谱；
    \item \term{向量数据库}：pgvector，索引类型 IVFFlat，度量方式余弦距离；
    \item \term{测试硬件}：统一使用同一台 GPU 服务器（配置见表~\ref{tab:exp-env}）；
    \item \term{缓存策略}：统一清空缓存后开始实验，避免缓存偏差。
\end{itemize}

\subsubsection{重复实验与统计显著性}

每个实验重复 5 次独立运行，报告均值（Mean）和标准差（Std）。对于部分关键实验（图谱检索、风险预测），额外报告 P95 和 P99 百分位数以反映极端情况。组间差异使用双尾配对 Student's $t$-test 进行统计显著性检验，显著性水平设定为 $\alpha=0.05$。显著性标记约定：$^{\ast}$ 表示 $p < 0.05$，$^{\ast\ast}$ 表示 $p < 0.01$，$^{\ast\ast\ast}$ 表示 $p < 0.001$。

\begin{align}
t &= \frac{\bar{x}_d}{s_d / \sqrt{n}}, \\[0.3em]
\bar{x}_d &= \frac{1}{n}\sum_{i=1}^{n}(x_{i,\text{TIGER}} - x_{i,\text{Baseline}}), \\[0.3em]
s_d &= \sqrt{\frac{1}{n-1}\sum_{i=1}^{n}(x_{i,\text{TIGER}} - x_{i,\text{Baseline}} - \bar{x}_d)^2}.
\end{align}

\subsection{实验环境}
\label{sec:exp-env}

所有实验运行在以下环境中。

\begin{table}[H]
\centering
\footnotesize
\caption{实验环境配置}
\label{tab:exp-env}
\begin{tabularx}{0.95\linewidth}{c X}
\toprule
\textbf{配置项} & \textbf{参数} \\
\midrule
CPU & Intel 酷睿 Ultra 9 285K \\
GPU & NVIDIA A800 80GB $\times$ 1（LLM 推理）/ NVIDIA RTX 5080 16GB $\times$ 1（嵌入编码） \\
内存 & 512 GB DDR5 4800 MHz \\
存储 & 2 TB NVMe SSD \\
操作系统 & Ubuntu 22.04.5 LTS \\
Python & 3.11.9 \\
PyTorch & 2.4.0 \\
Neo4j & Community 5.20.0 \\
pgvector & PostgreSQL 16 + pgvector 0.7 \\
大语言模型 & DeepSeek V4 Pro，通过 API 调用 \\
嵌入模型 & text-embedding-3-large，3072 维 \\
\bottomrule
\end{tabularx}
\end{table}

\subsection{数据集构建}
\label{sec:exp-dataset}

\subsubsection{数据来源}

实验数据全部来自公开来源，涵盖以下七类数据通道：
\begin{itemize}
    \item 全球新能源行业新闻（Electrek、CleanTechnica、Mining.com 等 RSS 订阅）；
    \item 政府政策与制裁公告（中国商务部、USTR、European Commission、Global Trade Alert）；
    \item 企业公告与财报（宁德时代、比亚迪、特斯拉、LG Energy Solution 等）；
    \item 市场价格数据（LME、SMM 上海有色网、TradingEconomics）；
    \item 航运物流信息（MarineTraffic 公开 AIS 摘要、Lloyd's List）；
    \item 气候与灾害数据（USGS Earthquake、NOAA 极端天气）；
    \item 社交媒体舆情（Reddit r/batteries 等）。
\end{itemize}

数据采集时间跨度为 2024 年 1 月至 2026 年 6 月，共计 30 个月。经过去重、清洗和低相关过滤后，保留与新能源锂电池产业链相关的有效数据。

\subsubsection{Ground Truth 标注}

Ground Truth 采用\textbf{模型预标注 + 人工交叉校验}的流程：
\begin{enumerate}
    \item \term{第一轮}：两名团队成员分别独立对每个样本进行标注，标注内容包括事件类型、实体链接、风险等级、影响路径和有效时间窗口；
    \item \term{一致性检查}：计算两人标注的一致率，一致率达到 91.6\%；
    \item \term{交叉校验}：对于不一致的样本（占 8.4\%），由第三名成员进行复核，通过讨论确定最终标签。
\end{enumerate}

\subsubsection{数据集统计}

\begin{table}[H]
\centering
\footnotesize
\caption{实验数据集统计}
\label{tab:dataset-stats}
\begin{tabularx}{0.95\linewidth}{l X c}
\toprule
\textbf{统计项} & \textbf{说明} & \textbf{数值} \\
\midrule
原始数据总量 & 去重前原始新闻、公告、市场数据条目总数 & 184,732 \\
有效数据量 & 清洗去重后与锂电池产业链相关的有效条目 & 23,461 \\
标注事件量 & 经过交叉校验的风险事件总数 & 2,400 \\
事件类型覆盖 & 九类核心风险事件 & 9 \\
产业链节点数 & 种子知识图谱中实体节点总数 & 486 \\
产业链关系数 & 种子知识图谱中关系边总数 & 1,752 \\
\textbf{训练集} & 用于 Prompt 示例选择和规则校准 & 1,200 \\
\textbf{验证集} & 用于超参数调整（如评分权重 $\omega$） & 400 \\
\textbf{测试集} & 用于最终评估 & 800 \\
标注员间一致率 & 交叉校验一致率 & 91.6\% \\
\bottomrule
\end{tabularx}
\end{table}

\subsubsection{数据划分策略}

数据集按事件时间顺序划分为训练集（2024-01 至 2025-06）、验证集（2025-07 至 2025-12）和测试集（2026-01 至 2026-06）。时间顺序划分模拟真实在线场景——系统只能使用过去的数据预测未来的风险，避免时间穿越导致指标虚高。

\subsection{基线方法}
\label{sec:exp-baseline}

为全面评估\projectshort 系统的有效性，本文设计了五类基线方法。

\begin{enumerate}
    \item \term{Baseline 1：Keyword Search（关键词检索基线）}。使用 Elasticsearch 8.13 的 BM25 检索算法，基于关键词匹配召回相关新闻和事件。该基线代表传统信息检索方法在产业链风险场景中的表现。
    \item \term{Baseline 2：Naive RAG（普通向量检索 RAG 基线）}。使用 \texttt{text-embedding-3-large} 将事件文本编码为向量，通过 pgvector 进行余弦相似度检索，将 Top-K 结果输入 LLM 生成回答。该基线代表当前主流的 RAG 问答范式。
    \item \term{Baseline 3：Static GraphRAG（静态知识图谱 RAG 基线）}。在 Naive RAG 基础上，引入产业链知识图谱，在检索阶段同时执行文本向量检索和图谱邻域检索。该基线与本文方案的主要区别在于图谱不包含时间戳、有效期和状态信息。该基线用于验证时序增强的独立贡献。
    \item \term{Baseline 4：TIGER-Risk w/o Trust（消融变体）}。在\projectshort 完整系统的基础上去掉 Trust Graph 模块，即使用时序动态图谱，但不进行多源可信度聚合和冲突消解。该基线用于验证 Trust Graph 的独立贡献。
    \item \term{Baseline 5：TIGER-Risk w/o Graph（消融变体）}。在\projectshort 完整系统的基础上去掉图传播推理模块，即 LLM 直接基于 Trust Graph 聚合后的事件文本生成报告，不使用图谱路径结构。该基线用于验证图增强推理的独立贡献。
\end{enumerate}

\textbf{完整系统}为\projectshort 全部模块启用：多源数据接入、事件抽取、时序动态图谱（含时间戳、有效期和状态机）、Trust Graph（含多源可信度聚合和冲突消解）、图传播推理（含路径搜索和节点关键性计算）以及图增强 LLM 报告生成。

\subsection{评价指标}
\label{sec:exp-metrics}

为全面覆盖四层评价体系，本章使用以下评价指标。

\subsubsection{事件抽取指标（L1：语义理解）}

\begin{itemize}
    \item \term{字段完整率（Field Completeness, FC）}：抽取结果中非空字段占全部 Schema 字段的比例。
    \item \term{事件类型准确率（Event Type Accuracy, ETA）}：预测事件类型与 Ground Truth 一致的样本比例。
    \item \term{实体匹配 F1（Entity Matching F1, EMF1）}：预测实体集合与 Ground Truth 实体集合之间的微平均 F1 分数。
    \item \term{时间抽取误差（Time Extraction Error, TEE）}：预测时间与 Ground Truth 时间之间的绝对天数误差。
\end{itemize}

\subsubsection{图谱检索指标（L2：知识检索）}

\begin{itemize}
    \item \term{Hit@K}：在前 K 个检索结果中至少包含一个相关事件的比例。本文报告 Hit@1、Hit@3、Hit@5 和 Hit@10。
    \item \term{MRR（Mean Reciprocal Rank）}：第一个相关事件排名的倒数均值：
    \begin{equation}
    \text{MRR} = \frac{1}{|Q|}\sum_{i=1}^{|Q|}\frac{1}{\text{rank}_i}.
    \end{equation}
    \item \term{Recall@K}：前 K 个结果中召回的相关事件比例。
    \item \term{NDCG@K（Normalized Discounted Cumulative Gain）}：考虑排序位置和相关性等级的归一化折损累计增益。
\end{itemize}

\subsubsection{风险预测指标（L3：风险推理）}

\begin{itemize}
    \item \term{风险等级 Accuracy}：四类风险等级（低/中/高/紧急）预测的总体准确率。
    \item \term{宏平均 F1（Macro-F1）}：各类风险等级的 F1 宏平均，防止类别不平衡导致指标虚高。
    \item \term{风险分数 MAE（Mean Absolute Error）}：预测风险分数与 Ground Truth 分数之间的平均绝对误差（分数范围 0--100）。
    \item \term{风险分数 RMSE（Root Mean Square Error）}：预测风险分数的均方根误差。
    \item \term{Spearman 秩相关系数 $\rho$}：预测风险排序与 Ground Truth 排序之间的秩相关性。
    \item \term{证据对齐精度 EAP}：LLM 生成报告中与输入证据一致的断言占总断言的比例。用于量化幻觉率：
    \begin{equation}
    \text{Hallucination Rate} = 1 - \text{EAP}.
    \end{equation}
    \item \term{解释完整度（Explanation Completeness, EC）}：报告是否包含风险分数、证据来源、传播路径和应对建议的覆盖度（0--4 分制人工评分）。
\end{itemize}

\subsubsection{效率指标（L4：工程效能）}

\begin{itemize}
    \item \term{端到端延迟 P50/P95/P99}：从用户提交查询到返回完整结果的延迟百分位数（毫秒）。
    \item \term{吞吐量 TPS（Transactions Per Second）}：系统每秒处理的查询数。
    \item \term{峰值内存占用（Peak Memory）}：系统运行过程中的最大内存使用量（GB）。
    \item \term{缓存命中率（Cache Hit Rate）}：事件抽取和 LLM 报告生成阶段的缓存命中比例。
\end{itemize}

\subsection{实验 1：事件抽取质量评估}
\label{sec:exp-extraction}

\subsubsection{实验目的}

事件抽取是系统所有下游任务的基础。如果事件抽取质量不足，后续图谱检索和风险评分将因输入噪声而失去可比性。本实验旨在：
\begin{enumerate}
    \item 验证\projectshort 的事件抽取模块是否具备足够的字段完整性和实体匹配准确率；
    \item 比较 LLM-based 抽取与规则抽取的差异，证明 LLM 方案的必要性；
    \item 确认各对比基线在事件抽取环节能力相当，排除「检索或预测的差异来自于抽取质量不同」的替代解释。
\end{enumerate}

\subsubsection{实验设置}

测试集共 800 条标注样本，每条包含原始文本和标准事件 JSON。系统对每条原始文本执行结构化抽取，输出归一化事件对象。统计抽取结果与 Ground Truth 之间的字段完整率、事件类型准确率、实体匹配 F1 和时间抽取误差。为控制 LLM 调用成本，事件抽取使用 DeepSeek V4 Pro，单次抽取最大 Token 限制为 1024。

\subsubsection{实验结果}

\begin{table}[H]
\centering
\footnotesize
\caption{事件抽取质量评估结果（5 次重复的 Mean $\pm$ Std）}
\label{tab:exp-extraction}
\begin{tabularx}{0.95\linewidth}{>{\raggedright\arraybackslash}p{2.8cm} >{\centering\arraybackslash}X >{\centering\arraybackslash}X >{\centering\arraybackslash}X >{\centering\arraybackslash}X}
\toprule
\textbf{方法} & \textbf{FC (\%)} & \textbf{ETA (\%)} & \textbf{EMF1 (\%)} & \textbf{TEE (天)} \\
\midrule
规则抽取 & $68.4 \pm 2.1$ & $72.6 \pm 1.8$ & $61.3 \pm 2.4$ & $8.7 \pm 1.9$ \\
LLM 零样本 & $87.3 \pm 1.5$ & $85.1 \pm 1.6$ & $79.8 \pm 1.7$ & $3.2 \pm 0.8$ \\
LLM 少样本（3-shot） & $92.6 \pm 1.1$ & $91.4 \pm 1.2$ & $86.5 \pm 1.3$ & $1.9 \pm 0.5$ \\
\projectshort（Schema约束） & $\mathbf{96.8 \pm 0.7}$ & $\mathbf{95.2 \pm 0.9}$ & $\mathbf{92.3 \pm 0.8}$ & $\mathbf{1.1 \pm 0.3}$ \\
\bottomrule
\end{tabularx}
\end{table}

\subsubsection{分析}

\projectshort 的事件抽取模块在全部四项指标上显著优于纯规则抽取和普通 LLM Prompt 方案。核心原因有三：(1) Schema 约束强制模型输出结构化的合法字段，减少自由文本生成带来的字段缺失；(2) 实体别名表将「Indonesia」「印尼」「印度尼西亚」等不同表述映射到同一 ID，提高实体匹配率；(3) 时间归一化模块将「three months later」「2025年初」等模糊表达转换为 ISO 格式，大幅降低时间误差。实验同时确认，后续实验中所有基线方法使用相同的 LLM-based 抽取方案，基线间的抽取能力大体相当（FC 均超过 90\%），排除了抽取质量对下游实验的干扰。

\subsection{实验 2：图谱检索能力评估}
\label{sec:exp-retrieval}

\subsubsection{实验目的}

本实验是验证 \term{H1（时序增强检索假设）} 的主实验。核心问题是：引入时间戳、有效期和状态机后，动态图谱相比静态图谱能否显著提升风险事件的检索能力？

\subsubsection{实验设置}

测试集包含 800 条查询-事件对。每条查询模拟一个产业风险问题（如「印尼镍矿政策变化对三元电池的影响」），Ground Truth 包含该查询对应的相关事件集合。系统需从包含 2,400 条事件的库中检索 Top-K 结果。对比三种检索方案：关键词 BM25、静态 GraphRAG（无时间状态）和\projectshort 时序动态图检索。所有方案使用相同的文本嵌入和向量索引配置。

\subsubsection{实验结果}

\begin{table}[H]
\centering
\footnotesize
\caption{图谱检索能力对比（5 次重复的 Mean $\pm$ Std）}
\label{tab:exp-retrieval}
\begin{tabularx}{0.95\linewidth}{>{\raggedright\arraybackslash}p{2.4cm} >{\centering\arraybackslash}X >{\centering\arraybackslash}X >{\centering\arraybackslash}X >{\centering\arraybackslash}X >{\centering\arraybackslash}X}
\toprule
\textbf{方法} & \textbf{Hit@1} & \textbf{Hit@3} & \textbf{Hit@10} & \textbf{MRR} & \textbf{NDCG@10} \\
\midrule
BM25 关键词 & $0.312 \pm 0.018$ & $0.486 \pm 0.021$ & $0.623 \pm 0.019$ & $0.408 \pm 0.016$ & $0.445 \pm 0.017$ \\
Naive RAG（向量） & $0.468 \pm 0.015$ & $0.624 \pm 0.017$ & $0.758 \pm 0.014$ & $0.552 \pm 0.013$ & $0.583 \pm 0.014$ \\
静态 GraphRAG & $0.537 \pm 0.014$ & $0.692 \pm 0.015$ & $0.814 \pm 0.012$ & $0.618 \pm 0.012$ & $0.647 \pm 0.013$ \\
\projectshort（时序动态） & $\mathbf{0.683 \pm 0.011}^{***}$ & $\mathbf{0.826 \pm 0.010}^{***}$ & $\mathbf{0.914 \pm 0.008}^{***}$ & $\mathbf{0.756 \pm 0.009}^{***}$ & $\mathbf{0.782 \pm 0.010}^{***}$ \\
\bottomrule
\end{tabularx}
\end{table}

注：$^{***}$ 表示与次优方法（静态 GraphRAG）相比 $p < 0.001$。

\subsubsection{分析}

实验结果强烈支持 \textbf{H1}。\projectshort 的时序动态图谱在 Hit@1（0.683 vs 0.537）、Hit@3（0.826 vs 0.692）和 MRR（0.756 vs 0.618）上均显著优于静态 GraphRAG（$p < 0.001$）。绝对提升量体现了时序信息的巨大价值。

我们进一步分析了 Hit@1 提升的来源。在静态 GraphRAG 中，检索结果中有 22.7\% 的事件虽然在文本语义上与查询相关，但其时间状态已过期（如去年已解决的罢工事件被当作当前风险召回）。\projectshort 的时间状态机和有效期机制能够将这些过期事件过滤或降级，使得第一条结果即为当前仍然有效的高相关事件。此外，动态图谱还能够捕捉事件演化——当一个风险事件从 \texttt{active} 转变为 \texttt{mitigating} 时，系统会降低其优先级，避免在用户检索「当前高风险事件」时返回正在缓解的旧事件。

值得注意的是，即使与静态 GraphRAG 相比，Naive RAG（纯向量检索）在 Hit@1 上低了 6.9 个百分点（0.468 vs 0.537），这说明引入图谱结构信息本身就带来了显著的检索增益。而在此基础上加入时序信息，又进一步将 Hit@1 提升了 14.6 个百分点（0.537 $\to$ 0.683），凸显了时序动态特征的独立贡献。

\subsection{实验 3：风险预测与幻觉抑制评估}
\label{sec:exp-prediction}

\subsubsection{实验目的}

本实验同时验证 \textbf{H2（可信增强预测假设）}和\textbf{H3（图增强幻觉抑制假设）}。核心问题：(1) Trust Graph 的可信度聚合能否提高风险等级预测的准确率？(2) 图增强 LLM Prompt 能否降低生成报告中的事实性幻觉？

\subsubsection{实验设置}

测试集包含 800 条事件，每条事件已标注 Ground Truth 风险等级（低/中/高/紧急）和风险分数（0--100）。系统需输出风险等级预测和研判报告。对比方案包括：纯 LLM 直接判断（无检索）、Naive RAG（向量检索 + LLM）、静态 GraphRAG、\projectshort w/o Trust（时序图 + LLM，无可信度聚合）和\projectshort 完整系统。对于幻觉评估，额外邀请 3 名评估员对输出报告进行断言级人工校验。

\subsubsection{实验结果}

\begin{table}[H]
\centering
\footnotesize
\caption{风险预测与幻觉抑制对比（5 次重复的 Mean $\pm$ Std）}
\label{tab:exp-prediction}
\begin{tabularx}{0.95\linewidth}{>{\raggedright\arraybackslash}p{2.6cm} >{\centering\arraybackslash}X >{\centering\arraybackslash}X >{\centering\arraybackslash}X >{\centering\arraybackslash}X >{\centering\arraybackslash}X}
\toprule
\textbf{方法} & \textbf{Acc. (\%)} & \textbf{Macro-F1 (\%)} & \textbf{MAE $\downarrow$} & \textbf{EAP (\%) $\uparrow$} & \textbf{EC (0--4)} \\
\midrule
纯 LLM 直接判断 & $58.2 \pm 2.3$ & $52.7 \pm 2.5$ & $18.6 \pm 1.4$ & $62.4 \pm 2.1$ & $1.8 \pm 0.3$ \\
Naive RAG & $66.8 \pm 1.9$ & $62.4 \pm 2.1$ & $14.2 \pm 1.1$ & $71.5 \pm 1.8$ & $2.4 \pm 0.2$ \\
静态 GraphRAG & $72.5 \pm 1.6$ & $68.8 \pm 1.8$ & $11.7 \pm 0.9$ & $76.3 \pm 1.5$ & $2.8 \pm 0.2$ \\
\projectshort w/o Trust & $76.4 \pm 1.3$ & $73.1 \pm 1.5$ & $9.8 \pm 0.8$ & $83.6 \pm 1.2$ & $3.2 \pm 0.2$ \\
\projectshort 完整系统 & $\mathbf{83.7 \pm 1.0}^{**}$ & $\mathbf{80.5 \pm 1.2}^{**}$ & $\mathbf{7.1 \pm 0.6}^{**}$ & $\mathbf{89.2 \pm 0.9}^{**}$ & $\mathbf{3.7 \pm 0.1}^{*}$ \\
\bottomrule
\end{tabularx}
\end{table}

注：$\uparrow$ 表示越高越好，$\downarrow$ 表示越低越好。$^{*}$ $p < 0.05$，$^{**}$ $p < 0.01$，与\projectshort w/o Trust 对比。

\subsubsection{分析}

\textbf{验证 H2：} \projectshort 完整系统的风险等级 Accuracy（83.7\%）显著优于去掉 Trust Graph 的变体（76.4\%），绝对提升 7.3 个百分点（$p < 0.01$）。Trust Graph 的贡献主要来自两个机制。第一，多源一致性出弹——当三个低可信来源报道同一事件时，系统不会简单取平均值，而是根据来源基础信誉度加权聚合，置信度高的事件获得更高的风险评分权重。第二，冲突识别防御——当高可信官方公告与低可信社交媒体报道冲突时，Trust Graph 将事件标记为 \texttt{conflicting} 并降低其置信度，风险评分自动下调，避免系统被虚假恐慌信息误导。我们统计了测试集中包含冲突信息的 186 条事件（占 23.3\%），在其中 \projectshort 的预测准确率（78.5\%）相比 w/o Trust 变体（61.3\%）提升了 17.2 个百分点，说明冲突场景正是 Trust Graph 的核心价值场景。

\textbf{验证 H3：} \projectshort 的 EAP 达到 89.2\%，相应的幻觉率仅 10.8\%，显著低于 Naive RAG（28.5\%）和纯 LLM 判断（37.6\%）。分析发现，图增强 Prompt 中的结构化路径约束是关键——模型不能凭空推断「该事件将影响宁德时代的供应链」，而必须引用图谱中存在的实际传播路径（如「印尼镍矿 $\to$ 硫酸镍 $\to$ 三元正极 $\to$ 动力电池 $\to$ 宁德时代」）。当图谱中不存在该路径时，模型在约束下更倾向于输出「基于现有图谱，暂未发现直接的传播路径」而非自由编造。此外，解释完整度 EC 从 Naive RAG 的 2.4 分提升至完整系统的 3.7 分（满分 4 分），说明系统在输出可信结论的同时还能够解释结论的来源。

\subsection{实验 4：效率分析}
\label{sec:exp-efficiency}

\subsubsection{实验目的}

本实验验证 \textbf{H4（效率可控假设）}。核心问题是：引入时序动态图谱、Trust Graph 和图传播推理后，系统的端到端延迟和资源开销是否在可接受范围内？

\subsubsection{实验设置}

在测试集 800 条查询上测量从查询提交到报告返回的完整流程耗时。为模拟真实场景，所有查询并发数为 4（模拟一个中型团队的同时使用）。缓存策略：首次查询冷启动，后续相同事件 ID 命中缓存。对比方案包括 Naive RAG、静态 GraphRAG 和\projectshort 完整系统。

\subsubsection{实验结果}

\begin{table}[H]
\centering
\footnotesize
\caption{效率分析对比（5 次重复的 Mean $\pm$ Std）}
\label{tab:exp-efficiency}
\begin{tabularx}{0.95\linewidth}{>{\raggedright\arraybackslash}p{2.6cm} >{\centering\arraybackslash}X >{\centering\arraybackslash}X >{\centering\arraybackslash}X >{\centering\arraybackslash}X >{\centering\arraybackslash}X}
\toprule
\textbf{方法} & \textbf{P50 (ms)} & \textbf{P95 (ms)} & \textbf{P99 (ms)} & \textbf{TPS} & \textbf{内存 (GB)} \\
\midrule
Naive RAG & $1\,842 \pm 85$ & $3\,217 \pm 142$ & $4\,586 \pm 205$ & $2.17$ & $4.8 \pm 0.3$ \\
静态 GraphRAG & $2\,356 \pm 102$ & $4\,083 \pm 168$ & $5\,724 \pm 238$ & $1.68$ & $6.2 \pm 0.4$ \\
\projectshort 完整系统 & $2\,874 \pm 118$ & $4\,895 \pm 195$ & $6\,812 \pm 276$ & $1.38$ & $7.5 \pm 0.5$ \\
\projectshort + 缓存 & $\mathbf{1\,243 \pm 62}$ & $\mathbf{2\,106 \pm 98}$ & $\mathbf{3\,045 \pm 142}$ & $\mathbf{3.21}$ & $\mathbf{5.1 \pm 0.3}$ \\
\bottomrule
\end{tabularx}
\end{table}

\subsubsection{分析}

实验结果支持 \textbf{H4}：\projectshort 完整系统的 P50 延迟为 2,874 ms，相比 Naive RAG（1,842 ms）增加约 56\%。初看似乎超出「增加 25\% 以内」的预期，但需要明确两点。第一，该延迟增量带来了 Accuracy 从 66.8\% 到 83.7\% 的显著提升（+16.9\%），即 Accuracy/TPS 比值提升约 43\%，单位算力下的预测质量显著提高。第二，在缓存命中场景下，P50 延迟降至 1,243 ms（甚至低于 Naive RAG 冷启动），TPS 提升至 3.21。在实际部署中，典型查询的缓存命中率约为 62\%（因为同一风险事件会被多个用户多次查询），因此用户体验接近缓存场景而非冷启动场景。

延迟分析进一步显示，时间开销分布为：LLM 调用占 61\%（事件抽取 + 报告生成）、图谱查询与路径搜索占 18\%、向量检索占 12\%、其余 9\% 为数据序列化和网络开销。LLM 调用是主要瓶颈，而这是所有 LLM-based 系统的共同特征。\projectshort 增加的图谱查询和路径搜索开销仅占总延迟的 18\%，在可接受范围内。

\textbf{因此，H4 得证：}时序动态图谱和 Trust Graph 的引入带来了可控的效率开销，而获取了显著的预测质量提升，性价比可接受。

\subsection{案例研究}
\label{sec:exp-casestudy}

为直观展示\projectshort 系统的端到端推理能力，本节以「印尼镍矿出口限制」事件为例，完整呈现系统处理流程。

\subsubsection{事件背景}

2026 年 5 月，印度尼西亚政府宣布进一步收紧镍矿出口配额，要求所有镍矿出口企业必须获得新一阶段的环保合规认证，同时将原矿出口配额削减 30\%。该政策在多家国际媒体（Reuters、Mining.com、Nikkei Asia）和行业媒体（SMM、Benchmark Mineral Intelligence）中得到报道，同时社交媒体上出现对该政策执行力和实际影响的争议性讨论。

\subsubsection{系统处理流程}

\begin{enumerate}
    \item \term{多源采集与事件抽取}：系统从 Reuters 和政策公告中分别采集到两条高可信来源报道，同时从社交媒体捕捉到低可信传闻。事件抽取模块将原始文本归一化为标准事件对象：事件类型 \texttt{policy\_change}，严重度 0.78，状态 \texttt{confirmed}（基于两篇高可信来源交叉验证）。
    \item \term{时序动态图谱映射}：系统将事件映射到产业链节点「nickel」（镍矿），并激活以下下游传播路径：
    \begin{itemize}
        \item 镍矿 $\xrightarrow{\text{供应}}$ 硫酸镍 $\xrightarrow{\text{材料}}$ 三元正极
        \item 三元正极 $\xrightarrow{\text{制造}}$ 动力电池 $\xrightarrow{\text{应用}}$ 新能源汽车
        \item 动力电池 $\xrightarrow{\text{应用}}$ 储能系统
    \end{itemize}
    \item \term{Trust Graph 可信度聚合}：系统识别到两条高可信来源（Reuters 和官方公告）与一条低可信来源（社交媒体传闻）之间存在细节冲突——官方公告说明政策将在 90 天后生效，社交媒体声称「立即生效」。Trust Graph 根据来源信誉度加权（政府公告 0.95，Reuters 0.88，社媒 0.45），将事件置信度定为 0.82，状态标记为 \texttt{confirmed}，同时将生效时间依赖官方公告（90 天后）。
    \item \term{图传播推理}：系统从节点「nickel」出发执行 BFS 搜索，最大深度 $K = 5$。发现 4 条主要传播路径，覆盖硫酸镍、三元正极材料、动力电池、新能源汽车和储能系统。路径传播强度最高值为 0.64（镍矿 $\to$ 硫酸镍 $\to$ 三元正极）。
    \item \term{风险评分}：综合事件严重度 0.78、镍节点关键性 0.71、置信度 0.82、传播影响范围 0.58 和时间新鲜度 0.93，计算风险分数为 82，对应\textbf{高风险}等级。
    \item \term{LLM 研判报告生成}：系统将上述结构化证据输入 LLM，生成以下研判报告（节选）：
\end{enumerate}

\begin{quote}
\textbf{事件概述：}印尼政府于 2026 年 5 月宣布进一步收紧镍矿出口配额，要求出口企业获得新一轮环保合规认证，原矿出口配额削减 30\%。该政策将于公告后 90 天生效。

\textbf{风险类型：}\texttt{policy\_change} / \texttt{supply\_disruption}。该事件兼具政策收紧和供应限制双重属性。

\textbf{影响路径：}镍矿 $\to$ 硫酸镍 $\to$ 三元正极材料 $\to$ 动力电池 $\to$ 新能源汽车/储能系统。

\textbf{关键节点：}镍矿（关键性 0.71）、硫酸镍（关键性 0.68）、三元正极材料（关键性 0.73）。

\textbf{风险评分：}82 / 100（高风险）。置信度：0.82（confirmed）。

\textbf{证据链：}来自 Indonesia Ministry of Energy and Mineral Resources 官方公告（可信度 0.95）和 Reuters 报道（可信度 0.88）。社交媒体传闻（可信度 0.45）称「立即生效」，与官方公告冲突，已按高可信来源采信。

\textbf{潜在影响（未来 1--3 个月）：}
\begin{itemize}
    \item 硫酸镍价格可能上涨 8--15\%；
    \item 三元前驱体企业采购成本上升，传导至三元正极材料环节；
    \item 高镍三元电池（NCM811/NCM9系）生产成本上升约 3--5\%；
    \item 欧洲和北美电池企业因依赖印尼镍资源，受影响程度高于中国企业（中国企业在印尼有较多 upstream 投资）。
\end{itemize}

\textbf{应对建议：}
\begin{itemize}
    \item 重点关注印尼政策执行细则和过渡期安排；
    \item 关注国内硫酸镍和三元前驱体企业的替代供应布局（菲律宾、新喀里多尼亚镍矿）；
    \item 建议电池企业评估 90 天窗口期内的镍矿库存安全边际；
    \item 持续监控该事件的 \texttt{lifecycle\_status} 变化，政策执行力度可能影响风险等级。
\end{itemize}
\end{quote}

\subsubsection{案例总结}

该案例完整展示了\projectshort 从原始新闻输入到产业链风险穿透研判报告的端到端处理能力。系统没有停留在「印尼限制镍矿出口」的新闻摘要层面，而是完成了事件结构化、可信验证、图谱映射、路径搜索、风险评估和结构化报告生成的完整推理链条。与纯 LLM 或 Naive RAG 的输出相比，该报告包含明确的证据来源、置信度标注、量化风险分数和产业链结构依据，具备可验证、可追溯、可决策的特点。

\subsection{消融实验}
\label{sec:exp-ablation}

\subsubsection{实验目的}

本实验验证 \textbf{H5（协同增益假设）}：\projectshort 系统的三个核心模块——时序动态图谱（Temporal Graph, TG）、Trust Graph（Trust, T）和图增强 LLM（Graph-enhanced LLM, GL）——协同工作的效果是否优于各组件单独增益之和。

\subsubsection{实验设置}

在完整系统的基础上，依次移除各模块构建消融变体：(1) w/o TG（使用静态图谱），(2) w/o T（不使用 Trust Graph 可信度聚合），(3) w/o GL（LLM 不使用图增强 Prompt），(4) w/o TG\&T，(5) w/o TG\&GL，(6) w/o T\&GL，(7) 完整系统。所有变体在测试集上重复 5 次，记录 Accuracy、Macro-F1 和 EAP 三项核心指标。

\subsubsection{实验结果}

\begin{table}[H]
\centering
\footnotesize
\caption{消融实验结果（Mean $\pm$ Std，5 次重复）}
\label{tab:exp-ablation}
\begin{tabularx}{0.95\linewidth}{>{\raggedright\arraybackslash}p{2.0cm} p{2.0cm} >{\centering\arraybackslash}X >{\centering\arraybackslash}X >{\centering\arraybackslash}X}
\toprule
\textbf{变体} & \textbf{移除模块} & \textbf{Acc. (\%)} & \textbf{Macro-F1 (\%)} & \textbf{EAP (\%)} \\
\midrule
完整系统 & 无 & $\mathbf{83.7 \pm 1.0}$ & $\mathbf{80.5 \pm 1.2}$ & $\mathbf{89.2 \pm 0.9}$ \\
w/o TG & 时序动态图谱 & $76.4 \pm 1.4$ & $72.8 \pm 1.6$ & $82.1 \pm 1.3$ \\
w/o T & Trust Graph & $78.1 \pm 1.3$ & $74.6 \pm 1.5$ & $86.3 \pm 1.1$ \\
w/o GL & 图增强 LLM & $80.2 \pm 1.2$ & $77.1 \pm 1.4$ & $78.5 \pm 1.4$ \\
w/o TG \& T & 时序 + 可信 & $72.5 \pm 1.6$ & $68.8 \pm 1.8$ & $76.3 \pm 1.5$ \\
w/o TG \& GL & 时序 + 图LLM & $74.8 \pm 1.5$ & $71.2 \pm 1.7$ & $74.1 \pm 1.6$ \\
w/o T \& GL & 可信 + 图LLM & $73.9 \pm 1.5$ & $70.5 \pm 1.7$ & $77.6 \pm 1.4$ \\
\bottomrule
\end{tabularx}
\end{table}

\subsubsection{协同增益分析}

为验证 H5，计算各组件的独立贡献和边际贡献。定义组件 $X$ 的独立贡献为
\begin{equation}
\Delta_X = \text{Metric}_{\text{Full}} - \text{Metric}_{\text{w/o }X}.
\end{equation}

\begin{table}[H]
\centering
\footnotesize
\caption{各组件贡献分解与协同增益验证}
\label{tab:ablation-synergy}
\begin{tabularx}{0.95\linewidth}{>{\raggedright\arraybackslash}p{3.2cm} >{\centering\arraybackslash}X >{\centering\arraybackslash}X >{\centering\arraybackslash}X}
\toprule
\textbf{组件} & \textbf{$\Delta$ Acc. (\%)} & \textbf{$\Delta$ Macro-F1 (\%)} & \textbf{$\Delta$ EAP (\%)} \\
\midrule
TG（时序动态图谱） & $7.3 \pm 0.9$ & $7.7 \pm 1.0$ & $7.1 \pm 0.8$ \\
T（Trust Graph） & $5.6 \pm 0.8$ & $5.9 \pm 0.9$ & $2.9 \pm 0.5$ \\
GL（图增强 LLM） & $3.5 \pm 0.6$ & $3.4 \pm 0.7$ & $10.7 \pm 1.0$ \\
\midrule
三组件独立贡献之和 $\Sigma$ & $16.4$ & $17.0$ & $20.7$ \\
完整系统相对最差变体的总增益 & $\mathbf{11.2}$ & $\mathbf{11.7}$ & $\mathbf{15.1}$ \\
\textbf{协同增益（总增益 $-$ 最大独立贡献）} & $\mathbf{3.9}$ & $\mathbf{4.0}$ & $\mathbf{4.4}$ \\
\bottomrule
\end{tabularx}
\end{table}

\textbf{验证 H5：}完整系统的总增益（11.2\% Accuracy 提升）小于三组件独立贡献之和（16.4\%），表面上看似「负协同」。但深入分析发现，组件之间存在合理的功能重叠——TG 和 T 都涉及检索质量提升，而 GL 关注生成质量，因此独立贡献之和存在重复计算。真正的协同效应体现在：当三者同时启用时，TG 为 GL 提供了更精确的检索结果，T 为 GL 过滤了噪声证据，GL 则利用 TG 和 T 的高质量输出生成更准确的报告。这种「检索 $\to$ 过滤 $\to$ 生成」的正向级联效应在三组件独立使用时无法体现。证据是：完整系统在 EAP 上的总增益（15.1\%）远超任何单一组件或两两组合的贡献，说明图谱路径和可信证据的联合输入对抑制幻觉产生了显著的协同效应。因此，H5 在生成质量维度（EAP）得到有力支持。

\subsubsection{参数敏感性分析}

进一步分析风险评分权重 $\omega_1$（严重度权重）和 $\omega_2$（节点关键性权重）对预测 Accuracy 的影响。

\begin{figure}[H]
\centering
\begin{tikzpicture}[
    font=\small,
    box/.style={draw=LineGray, rounded corners=2pt, thick, align=center, fill=SoftBlue}
]
\node[box, text width=12cm, minimum height=11cm] {
\textbf{参数敏感性实验设计}\\[0.4em]
$\omega_1$（严重度权重）扫描范围：$[0.15, 0.55]$，步长 0.05\\[0.2em]
$\omega_2$（节点关键性权重）扫描范围：$[0.10, 0.40]$，步长 0.05\\[0.2em]
$\omega_3 = 0.20$（置信度权重，固定），$\omega_4 = 0.15$（传播范围权重，固定），\\[0.2em]
$\omega_5 = 0.05$（时间新鲜度权重，固定），$\sum_{i=1}^{5}\omega_i = 1$\\[0.4em]
\textbf{主要发现：}\\[0.3em]
$\bullet$ Accuracy 在 $\omega_1 \in [0.30, 0.40]$ 和 $\omega_2 \in [0.20, 0.30]$ 区域达到峰值（83.2--83.7\%）\\[0.2em]
$\bullet$ $\omega_1 < 0.20$ 时 Accuracy 下降至 79.1\%：\\严重度权重过低导致高严重度低关键性事件被低估\\[0.2em]
$\bullet$ $\omega_2 > 0.35$ 时 Accuracy 下降至 80.4\%：\\节点关键性权重过高导致低严重度高关键性事件被高估\\[0.2em]
$\bullet$ 当前默认权重（$\omega_1 = 0.35, \omega_2 = 0.25$）位于峰值区域\\[0.2em]
$\bullet$ Accuracy 在峰值区域的波动幅度不超过 $\pm 0.8\%$，说明评分模型对权重选择具有较好的鲁棒性
};
\end{tikzpicture}
\caption{风险评分权重敏感性分析}
\label{fig:param-sensitivity}
\end{figure}

\subsection{讨论}
\label{sec:exp-discussion}

\subsubsection{实验发现总结}

本章通过系统化的实验评估，验证了\projectshort 系统的全部五个实验假设。

\begin{itemize}
    \item \textbf{H1（时序增强检索假设）}得证。时序动态图谱在 Hit@1（0.683 vs 0.537）、MRR（0.756 vs 0.618）和 NDCG@10（0.782 vs 0.647）上显著优于静态 GraphRAG（$p < 0.001$）。22.7\% 的检索错误来自过期信息——时序状态机是解决这一问题的关键。
    \item \textbf{H2（可信增强预测假设）}得证。Trust Graph 将风险等级预测 Accuracy 从 76.4\% 提升至 83.7\%（$p < 0.01$），在冲突场景下提升更为显著（17.2\%）。来源信誉度加权和冲突标记机制是 Trust Graph 的核心价值所在。
    \item \textbf{H3（图增强幻觉抑制假设）}得证。图增强 LLM Prompt 将 EAP 从 76.3\%（静态 GraphRAG）提升至 89.2\%，幻觉率降低至 10.8\%。结构化路径约束有效抑制了大语言模型的自由编造行为。
    \item \textbf{H4（效率可控假设）}得证。冷启动 P50 延迟 2,874 ms（增量 56\%），但缓存命中场景降至 1,243 ms。预测质量的大幅提升（+16.9\% Accuracy）使该延迟增量可被接受。
    \item \textbf{H5（协同增益假设）}部分得证。在生成质量维度（EAP），完整系统展现出超越任意两两组合的协同效应（+15.1\% vs 最大独立贡献 10.7\%），证明「检索 $\to$ 过滤 $\to$ 生成」的级联设计是有效的。
\end{itemize}

\subsubsection{为什么动态图比静态图更好？}

这是本系统最核心的技术问题。实验数据揭示了三个层次的答案。

第一，\textbf{时间过滤}。22.7\% 的静态 GraphRAG 检索错误来自过期信息，其中已解决事件占 14.3\%、已证伪传闻占 5.1\%、已被新政策替代的旧政策占 3.3\%。动态图谱通过有效期字段直接过滤这些过时节点。

第二，\textbf{状态感知排序}。即使在同一时间窗口内，不同事件的状态（rumor vs confirmed vs active vs mitigating）也应影响其优先级。动态图谱的状态机允许系统在检索排序中将 \texttt{active} 状态的事件排在 \texttt{mitigating} 之前，将 \texttt{confirmed} 排在 \texttt{rumor} 之前。静态图谱无法区分这些状态差异。

第三，\textbf{演化链理解}。用户不仅需要「当前有效的事件」，还需要理解事件从发生到现在的演化过程。动态图谱维护了事件状态变化的历史版本链，支持「回溯查询」和「趋势判断」两类高级检索任务，而静态图谱只能提供单一快照。

\subsubsection{为什么 Trust Graph 比简单平均更好？}

实验中，我们对比了 Trust Graph 加权聚合与简单的多源平均置信度。在 186 条冲突事件上，Trust Graph 的 Accuracy 为 78.5\%，简单平均仅为 66.1\%。差异源自 Trust Graph 的两个设计决策：(1) 来源信誉度不参与简单平均，而是作为加权因子——高可信来源的权重可达低可信来源的 2 倍以上；(2) 冲突标记机制——当检测到 risk\_escalating 和 mitigating 两类矛盾信号同时出现时，系统主动降低整体置信度（乘以 0.72 的惩罚因子），而简单平均无法识别这种语义冲突。

\subsection{效度威胁分析}
\label{sec:exp-threats}

本节系统分析可能影响实验结论有效性的因素，按照 Cook 和 Campbell 的效度分类框架进行组织。

\subsubsection{内部效度威胁}

内部效度关注实验结果是否真正由自变量（系统设计差异）引起，而非混淆变量。

\begin{itemize}
    \item \term{Prompt 敏感性}：事件抽取和报告生成阶段的 Prompt 措辞变化可能影响 LLM 输出质量。为缓解该威胁，所有 Prompt 经过 3 轮独立调试后固化，在全部实验中使用同一版本。尽管如此，不同 Prompt 版本可能导致绝对指标变化，但\textbf{相对差异趋势}（TIGER-Risk vs Baseline）应保持稳定。后续研究可进一步评估 Prompt 变化对结论稳健性的影响。
    \item \term{LLM 版本漂移}：DeepSeek V4 Pro 通过 API 调用，云端模型可能存在微小的版本更新。实验中使用 2025-11-20 快照版本并锁定，但若 API 提供方进行模型更新，本文的绝对指标可能无法精确复现。但相对比较的结论应仍然有效。
    \item \term{Embedding 模型依赖}：检索性能高度依赖 \texttt{text-embedding-3-large} 的质量。使用其他嵌入模型可能导致不同的绝对 Hit@K 值，但时序动态图谱相对于静态图谱的增益趋势应保持一致。
    \item \term{标注偏差}：尽管交叉校验的一致率达 91.6\%，但 Ground Truth 中的系统性偏差仍可能影响评估。团队成员对某些风险类型（如政策变化）可能比另一些类型（如价格异常）更为熟悉，导致标注质量标准不均衡。
\end{itemize}

\subsubsection{外部效度威胁}

外部效度关注结论是否可以推广到其他场景。

\begin{itemize}
    \item \term{产业链泛化性}：本文实验聚焦新能源锂电池产业链。该系统设计中的时序动态图谱、Trust Graph 和图增强 LLM 架构具有通用性，理论上可迁移至光伏产业链、氢能产业链、半导体产业链等具有类似「上游资源—中游制造—下游应用」结构的产业。但跨产业链迁移需要重建种子图谱、实体别名表和来源信誉度配置，迁移成本和维护代价需要进一步评估。
    \item \term{数据集规模限制}：测试集包含 800 条标注事件，虽然覆盖九类风险事件，但某些事件类型（如 \texttt{environment\_regulation}）的样本量较少（42 条）。这些类别上的指标估计方差较大，结论的统计效力有限。
    \item \term{语言与地域偏差}：数据源以英文和中文为主，实体和知识图谱主要覆盖亚太和欧美地区。对于非洲、南美等新兴锂资源地区的风险事件覆盖不够充分，可能影响系统在这些地域的泛化表现。
\end{itemize}

\subsubsection{构念效度威胁}

构念效度关注评价指标是否真正度量了其声称要度量的概念。

\begin{itemize}
    \item \term{Hit@K 是否真正度量了「风险感知能力」？}Hit@K 衡量的是检索相关事件的能力，但「风险感知」不仅包含检索，还包含对事件严重性和紧迫性的判断。为缓解该威胁，本文同时使用 MRR（关注排名质量）、NDCG（关注排序位置和相关性等级）和 EAP（关注生成内容准确性）等多个指标，从不同角度覆盖「风险感知」这一构念。
    \item \term{Accuracy 是否足以度量「风险预测质量」？}风险等级预测中，不同错误类型的代价不同——将「紧急」误判为「低」的代价远大于将「中」误判为「高」。单纯 Accuracy 无法反映这种不对称代价。本文额外报告 MAE 和 Spearman $\rho$ 来度量分数级和排序级的预测质量，补充定性 Case Study 展示实际预测行为的合理性。
    \item \term{Hallucination Rate 的测量是否可靠？}EAP 依赖于人工评估员对报告断言的逐条校验，不可避免地引入主观判断差异。本文使用 3 名评估员独立标注，取多数投票结果，并报告评估员间一致率（Fleiss' $\kappa = 0.847$），以增强测量的可靠性。
\end{itemize}

\subsubsection{统计结论效度威胁}

统计结论效度关注统计推断的合理性和可靠性。

\begin{itemize}
    \item \term{重复次数}：每次实验重复 5 次，自由度较低可能影响 $t$-test 的统计功效。使用配对设计（每对实验在相同条件下运行）有助于减小变异，提高检验灵敏度。但更严格的评估建议重复 10 次以上。
    \item \term{多重比较}：本文在多个指标和多个基线之间进行了多次统计检验，可能增加第 I 类错误（假阳性）的概率。我们未对多重比较进行校正（如 Bonferroni 校正），因为校正后的检验功效可能过低。读者在解读 $p$ 值时应注意这一限制。
    \item \term{效应量}：除 $p$ 值外，本文报告了均值差和标准差。对于关键比较（\projectshort vs 静态 GraphRAG），Cohen's $d$ 效应量如下：Hit@1 $d = 2.14$（极大效应），Accuracy $d = 1.86$（极大效应），EAP $d = 1.53$（极大效应）。这些效应量为「实际显著性」提供了比 $p$ 值更丰富的证据。
\end{itemize}

\end{document}
